\section{Related Work}
\label{sec:related}

\subsection{Test-time compute and self-consistency}

Scaling computation at inference time has become a standard tool for improving the reliability of large language models. Chain-of-thought prompting demonstrates that eliciting intermediate reasoning steps raises accuracy on arithmetic and symbolic tasks \citep{wei2022chain}, and self-consistency extends this observation by sampling multiple reasoning paths at nonzero temperature and returning the most frequent final answer \citep{wang2023selfconsistency}. Follow-up work reduces the number of samples needed through early stopping \citep{aggarwal2023lets}, generalizes aggregation to free-form generation by letting the model itself select among candidate answers \citep{chen2023universal}, and studies the optimal allocation of test-time compute relative to model size \citep{snell2024scaling}. Process reward models that score intermediate steps represent a further refinement of the same budget idea \citep{lightman2023lets}. Multi-agent debate aggregates answers across multiple model instances and rounds \citep{du2023improving}, and controlled comparisons find that debate does not reliably outperform self-consistency at matched budgets \citep{smit2024mad}, which reinforces the need for budget-matched comparisons when evaluating aggregation strategies. Our work differs from this literature in the source of diversity it studies: instead of asking how many samples to draw or how to weight them, we ask whether the same calls are better spent perturbing the prompt, and we answer under a matched-budget protocol.

\subsection{Prompt sensitivity and prompt ensembling}

A separate literature documents that model outputs are sensitive to semantically irrelevant details of the prompt. Reformatting alone, such as changing spacing or option ordering, shifts accuracy by several points \citep{sclar2024quantifying}, and the measured performance of a model depends strongly on which of many equivalent prompts an evaluator happens to choose \citep{mizrahi2024state}. Linguistic properties of the prompt, such as lexical familiarity, partially predict its success \citep{leidinger2023language}, and adversarial prompt perturbations can degrade performance further \citep{zhu2023promptbench}. The natural constructive response is to aggregate over prompts. Ask Me Anything generates many question reformulations with a language model and combines the resulting pseudo-labels \citep{arora2023ask}, an idea that predates and parallels self-consistency but targets zero-shot classification rather than matched-budget reasoning. Evaluation studies propose averaging over prompt templates to stabilize benchmark scores \citep{mizrahi2024state}. What remains missing is a direct comparison that holds the budget, the model, and the question fixed and contrasts prompt-space aggregation with answer-space aggregation on the same instances, which is the gap our protocol fills.

\subsection{Uncertainty estimation from sampled behavior}

Between the two aggregation literatures sits a third that uses sampling not to improve accuracy but to estimate uncertainty. Verbalized confidence asks the model to state a probability \citep{tian2023just,xiong2024can}, and self-probing methods ask the model to evaluate its own candidate answers \citep{kadavath2022language}. Behavioral methods infer uncertainty from the dispersion of sampled outputs, from simple answer agreement to semantic clustering that treats paraphrased answers as equivalent \citep{kuhn2023semantic,farquhar2024detecting}, with surveys consolidating the rapidly growing set of measures \citep{lin2024uncertainty}. Calibration metrics such as expected calibration error \citep{guo2017calibration} and selective prediction metrics such as the area under the risk-coverage curve quantify how well a confidence score ranks correct above incorrect answers. These methods inherit an assumption from sampling: that the sample distribution is informative. Our study tests this assumption directly, characterizes the regimes in which sample agreement and paraphrase agreement carry signal, and shows that the two agreement signals are complementary where both vary. In this sense our fragility analysis also connects to work on the faithfulness of chain-of-thought explanations, which finds that stated reasoning does not always cause the final answer \citep{lanham2023measuring}: paraphrase stability measures the robustness of the answer to surface form rather than to the reasoning trace.

\subsection{Positioning}

In summary, answer-space ensembling, prompt-space ensembling, and sampling-based uncertainty estimation have developed largely in isolation. Answer-space methods assume resampling generates useful diversity. Prompt-space methods assume paraphrases probe the same knowledge. Uncertainty methods assume sample dispersion tracks correctness. All three assumptions are regime dependent, and no prior work, to our knowledge, measures all three under a single matched-budget protocol on modern models. The comparison also carries practical weight for evaluation practice, since benchmark conclusions about model capability shift when prompts vary \citep{mizrahi2024state} and when aggregation is applied \citep{wang2023selfconsistency}. Our contribution is the unified empirical map, the validation machinery for paraphrases, and the decision guidance that follows from the regime structure.
